% =============================================================================
% CAPÍTULO 2 — SPRINT 1: Cimientos, Seguridad y Flujos Centrales
% =============================================================================
\chapter{Sprint 1 — Cimientos, Seguridad y Flujos Centrales}
\label{cap:sprint1}
\resumenCapitulo{Primer incremento funcional: gestión de cuentas de usuario y perfil personal,
sustentada en la base de autenticación, la seguridad JWT y los flujos centrales de acceso de
la plataforma.}

El Sprint 1 entrega el primer incremento de software con valor para el usuario: el sistema de
\textbf{identidad, acceso y seguridad} (EPIC-01), cimiento sobre el que se construyen todos los
flujos posteriores. Se materializa la autenticación delegada en Supabase, la validación de JWT
en el backend como \textit{Resource Server} y los layouts base del frontend.\par

\textbf{Periodo de ejecución:} del lunes 30 de marzo al sábado 11 de abril de 2026 (2 semanas).
\textbf{Módulo entregado (Plan General):} \textit{Cuentas de Usuario y Perfil Personal}.\par

\section{Resumen del Incremento y Funcionalidades Abordadas}
\label{sec:s1_resumen}

Durante este incremento (\textit{tag} \comando{release/v1.0.0}) se completaron las Historias
LOG-001 a LOG-011: registro de profesionales y reclutadores, inicio de sesión, autenticación con
proveedores externos (OAuth), recuperación de contraseña, cierre de sesión, control de acceso por
rol, cambio de contraseña/email con OTP y eliminación de cuenta. La tabla~\ref{tab:s1_alcance}
resume el alcance abordado por capa, que se desarrolla en las
Secciones~\ref{sec:s1_database}--\ref{sec:s1_gestion}.

\begin{table}[!ht]
    \centering
    \renewcommand{\arraystretch}{1.35}
    \fontsize{9}{10.8}\selectfont\sffamily
    \begin{tabularx}{\textwidth}{|l|X|}
        \hline
        \rowcolor{colorPrimario}
        \textcolor{white}{\textbf{Capa}} & \textcolor{white}{\textbf{Funcionalidad entregada}} \\ \hline
        Base de datos & Esquema de roles y políticas RLS para datos sensibles (Sección~\ref{sec:s1_database}). \\ \hline
        \rowcolor{grisTecnico}
        Backend & API de \comando{auth} y \comando{security}; filtros JWT y CORS (Sección~\ref{sec:s1_backend}). \\ \hline
        Frontend & Navbar, Sidebar responsiva, formularios de auth y de perfil (Sección~\ref{sec:s1_frontend}). \\ \hline
        \rowcolor{grisTecnico}
        ISO 26515 & Borrador de la spec OpenAPI y diagrama de componentes de seguridad (Sección~\ref{sec:s1_gestion}). \\ \hline
    \end{tabularx}
    \caption{Alcance del Sprint 1 por capa.}
    \label{tab:s1_alcance}
\end{table}

\subsection{Objetivos del Incremento}
\label{ssec:s1_objetivos}
Los objetivos concretos comprometidos para este sprint fueron:

\begin{enumerate}[noitemsep]
    \item \textbf{Identidad federada:} integrar Supabase Auth como proveedor de identidad y
    configurar el backend como \textit{Resource Server} que valida JWT vía JWKS.
    \item \textbf{Registro diferenciado por rol:} habilitar el alta de cuentas
    \comando{PROFESSIONAL} y \comando{RECRUITER} con redirección posterior según el rol.
    \item \textbf{Operaciones sensibles protegidas:} implementar el flujo OTP para cambio de
    contraseña, cambio de email y eliminación de cuenta.
    \item \textbf{Blindaje perimetral:} desplegar la cadena de filtros de seguridad (cookie de
    token, cabeceras, no-cache, CORS) y la política RBAC de administración.
    \item \textbf{Cimiento de la interfaz:} entregar los layouts base (Navbar y Sidebar
    responsiva) y los formularios de autenticación y perfil.
\end{enumerate}

\clearpage
\subsection{Historias de Usuario Abordadas}
\label{ssec:s1_backlog}
La tabla siguiente recoge las Historias de la EPIC-01 comprometidas y cerradas en el sprint,
con su esfuerzo en puntos y criterio de aceptación.

\begingroup
\footnotesize
\renewcommand{\arraystretch}{1.35}
\begin{TablaBacklog}
    LOG-001 & Registro de usuario profesional & Alt & 5 & 1 & Creación exitosa tras validación. \\ \hline
    LOG-002 & Registro de usuario reclutador & Alt & 5 & 1 & Registro con rol asignado. \\ \hline
    LOG-003 & Inicio de sesión & Alt & 3 & 1 & Acceso con credenciales válidas. \\ \hline
    LOG-004 & Autenticación con proveedores ext. & Med & 5 & 1 & Login vía OAuth exitoso. \\ \hline
    LOG-005 & Recuperación de contraseña & Alt & 8 & 1 & Cambio de clave por enlace seguro. \\ \hline
    LOG-006 & Cierre de sesión & Alt & 2 & 1 & Invalida token (blacklist) y redirige. \\ \hline
    LOG-007 & Asignación de roles & Alt & 8 & 1 & Rol según tabla de perfil (basic/company). \\ \hline
    LOG-008 & Control de acceso por rol & Alt & 8 & 1 & Bloqueo de rutas sin permisos. \\ \hline
    LOG-009 & Cambio de contraseña/email (OTP) & Alt & 8 & 1 & Verificación por código OTP. \\ \hline
    LOG-010 & Eliminación de cuenta & Med & 8 & 1 & Baja confirmada con OTP. \\ \hline
    LOG-011 & Protección de credenciales & Alt & 8 & 1 & Cifrado delegado en Supabase. \\ \hline
\end{TablaBacklog}
\endgroup

\begin{NotaTecnica}
La velocidad comprometida (60 puntos de historia) se ajustó a la capacidad real del equipo tras
la incepción. El criterio de cierre exige que cada Historia satisfaga la \textit{Definition of
Done} transversal fijada en la Sección~\ref{ssec:s0_dod}.
\end{NotaTecnica}

\subsection{Ceremonias del Sprint}
\label{ssec:s1_ceremonias}
La tabla~\ref{tab:s1_ceremonias} resume las ceremonias Scrum del incremento (30 mar--11 abr 2026).

\begin{table}[!ht]
    \centering
    \renewcommand{\arraystretch}{1.3}
    \fontsize{9}{10.8}\selectfont\sffamily
    \begin{tabularx}{\textwidth}{|l|X|}
        \hline
        \rowcolor{colorPrimario}
        \textcolor{white}{\textbf{Ceremonia}} & \textcolor{white}{\textbf{Resultado}} \\ \hline
        Planning & Selección de las Historias LOG-001--011 y estimación. \\ \hline
        \rowcolor{grisTecnico}
        Daily & Sincronización diaria; foco en la integración de Supabase Auth. \\ \hline
        Review & Demostración del registro, login y OTP al Product Owner. \\ \hline
        \rowcolor{grisTecnico}
        Retrospective & Acuerdo sobre convenciones de seguridad y filtros. \\ \hline
    \end{tabularx}
    \caption{Ceremonias Scrum del Sprint 1.}
    \label{tab:s1_ceremonias}
\end{table}

\clearpage
